%! Prerequisite Skills Review — Unit 1, Lesson 1.0 — Group Activity (Answer Key)
\documentclass[10pt]{article}
\usepackage{precalculus-article}
\usepackage{precalculus-key}   % pulls in -boxes; do NOT also load -boxes
\newcommand{\TallMath}[1]{$\displaystyle #1\rule[-1.4em]{0pt}{3.2em}$}

\begin{document}

\pageheader{Unit 1, Lesson 1.0}{Group Activity --- Answer Key}

\vspace{0.12in}

\begin{scenariobox}[Car-Wash Fundraiser]{sky}
\small
The math club runs a Saturday car wash. Supplies cost \$40, and the club charges
\$8 per car, so the club's profit after washing $c$ cars is
\[ P = 8c - 40. \]
Work through the tier your teacher assigns. Every tier uses this same formula.
\end{scenariobox}

\vspace{0.1in}

\boxguard
\begin{tcolorbox}[colback=white, colframe=black!40,
    title=\textbf{Tier R --- Remediate},
    fonttitle=\bfseries, arc=2mm, left=3mm, right=3mm, top=2mm, bottom=2mm]

\begin{enumerate}[itemsep=8pt]
  \item Complete the table of profits.

        \smallskip
        \renewcommand{\arraystretch}{1.4}
        \begin{tabular}{|c|c|c|c|c|}
        \hline
        Cars washed $c$ & 0 & 5 & 10 & 20 \\
        \hline
        Profit $P$ (\$) & \ans{$-40$} & \ans{0} & \ans{40} & \ans{120} \\
        \hline
        \end{tabular}

  \item The profit at $c = 0$ is negative. What does that number mean for the club?

        \vspace{0.05in}
        \ansline{If no cars are washed, the club loses the \$40 it spent on supplies.}

  \item Solve $8c - 40 = 0$ to find the \textbf{break-even point} --- the number of
        cars where the club stops losing money.

\begin{work}
  8c - 40 &= 0 \\
       8c &= 40 \\
        c &= 5
\end{work}

        Break-even: \ans{5} cars
\end{enumerate}
\end{tcolorbox}

\vspace{0.1in}

\boxguard
\begin{tcolorbox}[colback=white, colframe=black!40,
    title=\textbf{Tier A --- Approaching Proficiency},
    fonttitle=\bfseries, arc=2mm, left=3mm, right=3mm, top=2mm, bottom=2mm]

\begin{enumerate}[itemsep=8pt]
  \item The club's goal is a \$120 profit. Solve $8c - 40 = 120$ to find how many
        cars that takes.

\begin{work}
  8c - 40 &= 120 \\
       8c &= 160 \\
        c &= 20
\end{work}

        Cars needed: \ans{20}

  \item A rival goal: profit of \textbf{at least} \$200. Solve $8c - 40 \ge 200$ and
        write the solution in interval notation.

\begin{work}
  8c - 40 &\ge 200 \\
       8c &\ge 240 \\
        c &\ge 30
\end{work}

        Interval notation: \ans{$[30, \infty)$}

  \item Write one sentence interpreting your interval answer in terms of cars washed.

        \vspace{0.05in}
        \ansline{The club must wash 30 or more cars to make a profit of at least \$200.}
\end{enumerate}
\end{tcolorbox}

\vspace{0.1in}

\boxguard
\begin{tcolorbox}[colback=white, colframe=black!40,
    title=\textbf{Tier E --- Extension},
    fonttitle=\bfseries, arc=2mm, left=3mm, right=3mm, top=2mm, bottom=2mm]

The science club runs a rival car wash with cheaper supplies but a lower price:
their profit is $P = 6c + 10$.

\begin{enumerate}[itemsep=8pt]
  \item Solve $8c - 40 = 6c + 10$ to find the number of cars where the two clubs
        earn the same profit.

\begin{work}
  8c - 40 &= 6c + 10 \\
       2c &= 50 \\
        c &= 25
\end{work}

        Same profit at $c = $ \ans{25} cars; that profit is \$\ans{160}

  \item For more than that many cars, which club earns more? Explain using the two
        formulas --- not by testing single values.

        \vspace{0.05in}
        \ansline{The math club. Its profit grows \$8 per car versus \$6 per car, so once the clubs tie at 25 cars, the math club gains \$2 more with every additional car and stays ahead for all $c > 25$.}
\end{enumerate}
\end{tcolorbox}

\end{document}
